\subsection*{Direct comparison with causal sampling and sequential experiment order}
The density cap and conditional-hazard construction have direct causal-sampling antecedents, not only a connection to trimming. For i.i.d. proposals of law $\nu$ and a stopping index $K$ returning the current sample with law $\lambda$, the standard predictable-count argument gives
\[
 \lambda(A)\le\E K\,\nu(A),\qquad
 \E K\ge\|d\lambda/d\nu\|_\infty
             =2^{D_\infty(\lambda\|\nu)}.
\]
The last equality uses base-two divergence. This is the exact runtime lower bound in Goc--Flamich \cite[Theorem III.1]{GocFlamich24}; it is not new here. If $d\lambda/d\nu\le b$, the ordinary acceptance probability $(d\lambda/d\nu)/b$ realizes $\lambda$ at mean index $b$. Consequently the i.i.d. fixed-dictionary special case of $\Q_m^\ell(b)$ is optimization of the task loss over this classical feasible sampling set. Liu--Verd\'u \cite{LiuVerdu18} study moments of causal and noncausal selected indices and sampling--covering duality. The profile alone, its selection threshold and its i.i.d. density-cap converse are therefore not asserted as independent new principles.

The distinction in the transfer theorem is its full-current-report, actual-history comparison measure for an arbitrary competing controller under one actual operating model, coupled to a different common implementation under every model. This measure need not be the law of any executable reference stopping rule. Decoder cardinality, retained live labels and code length of a shared-randomness index are different resources. Goc--Flamich's channel-simulation divergence bounds index entropy; neither $m$ nor the total number of controller labels in this article is that entropy. Arbitrary noncausal selection of a past proposal is not admitted merely because its final output distribution agrees.

The report cites the first version of Li's binary-experiment preprint. Its checked successor is Li--Liu, \emph{Ranking Experiments under Sequential Sampling}, arXiv:2608.19897v2, dated 28 August 2026 \cite{LiLiu26}. Their Theorem 1 treats generic pairs of finite, full-support, pairwise identified \emph{multistate} experiments; it equates the stated large-budget stopping and small-cost decision orders with strict inequalities of every directed pairwise KL divergence. Theorem 2 gives exact reproduction of arbitrary finite-output stopping policies with a statewise KL-ratio leading cost and a uniform square-root remainder. Restricting that current comparison to binary experiments would be inaccurate.

That theorem ranks repeatable i.i.d. sources and reproduces entire terminal experiments. Our localized fixed-loss task order is not Blackwell order, Le Cam deficiency, their generic sequential order, or a complete quantum sufficiency criterion. Theorems~\ref{thm:width-defect} and~\ref{thm:quadratic-defect} identify exact defects for their specified tasks, while Theorem~\ref{thm:primitive-risk} gives a finite-budget approximation for controlled state-changing instruments with counted finite implementations. It does not prove a competing KL characterization or the full sequential-conversion theorem.

The invariant-subspace obstruction used for the new fallback is also established quantum reachability machinery \cite[Theorem 3.6]{Ying14pre}. The new statement is the resulting explicit same-budget risk sandwich with a derived fallback, a common policy across finitely many unknown models and a fixed output width. Its proof separates the positive survival mass bound from the generally unavailable uniform expected cost-tail bound. This specifies a mathematical assertion for independent priority and significance assessment, not a claim to originate compactness, sequence form or quantum reachability.
